\documentclass[12pt]{article}
\usepackage{amsmath}
\usepackage{amssymb}
\usepackage{amsfonts}
\usepackage{graphicx}
\usepackage{hyperref}
\usepackage{geometry}
\usepackage{esint}
\usepackage[portuguese]{babel}
\usepackage[utf8]{inputenc}

\geometry{
 a4paper,
 left=3cm,
 top=2cm,
 right=2cm,
 bottom=2cm
}

\begin{document}

\begin{center}
    \textbf{\large UNIVERSIDADE FEDERAL DO RIO GRANDE DO SUL} \\
    \textbf{\large INSTITUTO DE FÍSICA} \\
    \vspace{3cm}
    \textbf{\Large FÍSICA GERAL - ELETROMAGNETISMO} \\
    \vspace{1cm}
    \textbf{\large REVISÃO} \\
    \vspace{3cm}
    
    \begin{tabular}{c}
        LEONARDO KAWAI SUAREZ        
    \end{tabular}
    
    \vfill
    \textbf{\large Porto Alegre, RS} \\
    \textbf{\large\today} \\
\end{center}\newpage


\section*{Área 1}

\subsection*{Lei de Coulomb}
\subsection*{Campo Elétrico}
\subsection*{Lei de Gauss}

\newpage

\section*{Área 2}

\subsection*{Potencial Elétrico}
\subsection*{Capacitância}
\subsection*{Circuitos de Corrente Contínua}

\newpage

\section*{Área 3}

\subsection*{Campo Magnético}

A Lei de Biot-Savart descreve o campo magnético \(\mathbf{B}\) gerado por uma corrente elétrica \(I\) em um ponto \(P\) a uma distância \(\mathbf{r}\) do elemento de corrente \(d\mathbf{l}\):

\begin{equation}\label{eq:lei_biot_savart}
    \boxed{\mathbf{B}(r) = \frac{\mu_0}{4\pi} \int \frac{I d\mathbf{l} \times \mathbf{r}}{r^3}}
\end{equation}

\noindent onde \(\mu_0\) é a permeabilidade magnética do vácuo, \(d\mathbf{l}\) é o elemento de corrente e \(\mathbf{r}\) é o vetor que aponta do elemento de corrente ao ponto \(P\).

Podemos utilizar a Lei de Biot-Savart para calcular o campo magnético gerado por diversas configurações de corrente elétrica, como espiras circulares e fios retilíneos infinitos.

\begin{itemize}
    \item [(a)] \textbf{Campo magnético de uma espira circular}:
        \begin{equation}
            \boxed{\mathbf{B}(z) = \frac{\mu_0 I R^2}{2(R^2 + z^2)^{3/2}} \hat{\mathbf{z}}}
        \end{equation}
        \noindent onde \(R\) é o raio da espira, \(I\) é a corrente elétrica e \(z\) é a distância do ponto ao plano da espira.
    \item [(b)] \textbf{Campo magnético ao redor de um fio retilíneo infinito}:
        \begin{equation}
            \boxed{\mathbf{B}(r) = \frac{\mu_0 I}{2\pi r}\hat{\mathbf{\phi}}}
        \end{equation}
        \noindent onde \(r\) é a distância do ponto ao fio.
\end{itemize}

\subsection*{Lei de Ampère}

A Lei de Ampère é uma das equações de Maxwell e relaciona o campo magnético \(\mathbf{B}\) com a corrente elétrica \(I\) que atravessa uma curva fechada \(C\), comumente chamada de curva amperiana. A lei é expressa pela seguinte equação:

\begin{equation}
    \boxed{\oint_{C} \mathbf{B} \cdot d\mathbf{l} = \mu_0 I_{C}}
\end{equation}

\noindent onde \(I_{C}\) é a corrente elétrica que atravessa a superfície delimitada pela curva amperiana \(C\).

Em problemas envolvendo simetria, podemos aplicar a Lei de Ampère para calcular o campo magnético gerado por diferentes condutores, como fios retilíneos infinitos, toróides e solenoides.

\begin{itemize}
    \item [(a)] \textbf{Campo magnético de um fio retilíneo infinito}:\\
        \begin{align}
                \boxed{\mathbf{B}(r) = \begin{cases}
                    \frac{\mu_0 I}{2\pi} \frac{r}{R^2} \hat{\mathbf{\phi}} \quad &(0 \leq r \leq R) \\
                    \frac{\mu_0 I}{2\pi} \frac{1}{r} \hat{\mathbf{\phi}} \quad &(R \leq r)
                \end{cases}}
        \end{align}    

        \noindent onde \(R\) é o raio do condutor e \(r\) é a distância do ponto ao fio.

    \item [(b)] \textbf{Campo magnético de um toróide}:
        \begin{align}
            \boxed{\mathbf{B}(r) = \begin{cases}
            		\frac{\mu_0 N I}{2\pi r}\hat{\mathbf{\theta}} \quad &(a < r < b) \\
            		0 \quad &(r < a \,\text{ ou }\, r > b)
            	\end{cases}}
        \end{align}
        \noindent onde \(N\) é o número de espiras do toróide de raios interno \(a\) e externo \(b\) e \(r\) é a distância do ponto ao eixo do toróide.

    \item [(c)] \textbf{Campo magnético de um solenoide}:
        \begin{align}
            \boxed{\mathbf{B} = \begin{cases}
            		\mu_0 n I\hat{\mathbf{z}} \quad &(\text{dentro do solenoide}) \\
            		0 \quad &(\text{fora do solenoide})
            	\end{cases}}
        \end{align}
        \noindent onde \(n\) é o número de espiras por unidade de comprimento do solenoide.
\end{itemize}

A \emph{força por unidade de comprimento} entre dois condutores paralelos infinitos é dada por:

\begin{equation}
    \boxed{f = I_1 B_2 = \frac{\mu_0 I_1 I_2}{2\pi d}}
\end{equation}

\noindent onde \(I_1\) e \(I_2\) são as correntes nos condutores e \(d\) é a distância entre eles.

\newpage

\section*{Área 4}

\subsection*{Indução Eletromagnética}

A Lei da Indução de Faraday afirma que a força eletromotriz induzida em um circuito fechado é proporcional à taxa de variação do fluxo magnético através do circuito.

\begin{equation}\label{eq:lei_faraday}
    \boxed{\mathcal{E} = -\frac{d\Phi_{B}}{dt}}
\end{equation}

\noindent onde \(\mathcal{E}\) é a força eletromotriz induzida e \(\Phi_{B}\) é o fluxo magnético que atravessa o circuito. A unidade de fluxo magnético é o weber (Wb), que é equivalente a \(1 \, \text{Wb} \equiv 1 \, \text{T}\cdot\text{m}^2\).

O sinal negativo de \eqref{eq:lei_faraday} indica que a corrente induzida tem um sentido tal que o campo magnético gerado por ela se opõe à variação do fluxo magnético que a causou, conforme enunciado pela Lei de Lenz.

\subsection*{Indutância}

Indutor é um componente elétrico que armazena energia na forma de um campo magnético quando uma corrente elétrica passa por ele. A indutância \(L\) de um indutor é uma medida de sua capacidade de induzir uma força eletromotriz (fem) em resposta a uma variação da corrente elétrica.

A indutância é definida como

\begin{equation}
    \boxed{L \equiv \frac{N\Phi_{B}}{I}}
\end{equation}

\noindent onde \(N\) é o número de espiras, \(\Phi_{B}\) é o fluxo magnético que atravessa cada espira e \(I\) é a corrente elétrica. A unidade de indutância é o henry (H), que é equivalente a \(1 \, \text{H} \equiv 1 \, \text{Wb}/\text{A}\).

\subsection*{Circuitos de Corrente Alternada}



\subsection*{Equações de Maxwell}

\begin{itemize}
    \item [(a)] \textbf{Lei de Gauss para o campo elétrico}
        \begin{equation}
            \boxed{\oint_{S} \mathbf{E}\cdot d\mathbf{S} = \frac{Q_{int}}{\varepsilon_0}}
        \end{equation}
        \noindent onde \(Q_{int}\) é a carga interna à superfície \(S\).

    \item [(b)] \textbf{Lei de Gauss para o campo magnético}
        \begin{equation}
            \boxed{\oint_{S} \mathbf{B} \cdot d\mathbf{S} = 0}
        \end{equation}
        \noindent (não existem monopolos magnéticos).
    \item [(c)] \textbf{Lei de Faraday}
        \begin{equation}
            \boxed{\oint_{C} \mathbf{E} \cdot d\mathbf{l} = -\frac{d}{dt}\int_{S} \mathbf{B} \cdot d\mathbf{S}}
        \end{equation}
    \item [(d)] \textbf{Lei de Ampère-Maxwell}
        \begin{equation}
            \boxed{\oint_{C} \mathbf{B} \cdot d\mathbf{l} = \mu_0 \left( I_{cond} + I_{d} \right)}
        \end{equation}
        \noindent onde \(I_{d} = \varepsilon_0 \frac{d}{dt}\int_{S} \mathbf{E} \cdot d\mathbf{S}\) é a corrente de deslocamento.
\end{itemize}



\end{document}